\begin{ouverture}[calcul]{2}{Multiples, diviseurs et partages}{Partie 1 · Nombres et calculs}
{Division euclidienne · Critères · Recherche organisée}
\objectif{Utiliser quotient et reste.}
\objectif{Reconnaître des multiples et des diviseurs.}
\objectif{Organiser une recherche de solutions entières.}
\prerequis{Tables et calcul posé.}
\end{ouverture}
\pagestyle{otMain}\essentiel\fontsize{11}{15}\selectfont
\phantomsection\label{cours-2-1}
\essSection{1}{Comprendre la division entière}
\begin{definition}Pour des entiers $a\ge0$ et $b>0$, la division euclidienne donne $a=bq+r$, avec $q$ entier et $0\le r<b$. $q$ est le quotient, $r$ le reste.\end{definition}
\begin{exemple}$83=7\times11+6$. On fait 11 paquets de sept et il reste six objets. Si tous les objets doivent être transportés dans des boîtes de sept, il faut douze boîtes.\end{exemple}
\begin{definition}$a$ est un multiple de $b$ si $a=bq$ pour un entier $q$. Pour $b>0$, cela signifie que $b$ divise $a$ et que le reste de la division est nul. Ainsi 36 est multiple de 9, et 9 est un diviseur de 36.\end{definition}
\begin{savoirfaire}[Ne rien oublier]Chercher les diviseurs par paires : $24=1\times24=2\times12=3\times8=4\times6$. La liste complète est $1,2,3,4,6,8,12,24$.\end{savoirfaire}
\clearpage
\phantomsection\label{cours-2-2}
\essSection{2}{Critères et recherches}
\begin{propriete}Un entier est divisible par 2 si son chiffre des unités est pair ; par 5 s’il finit par 0 ou 5 ; par 10 s’il finit par 0. Il est divisible par 3, respectivement par 9, si la somme de ses chiffres est divisible par 3, respectivement par 9.\end{propriete}
\begin{exemple}Pour 234, $2+3+4=9$ : il est divisible par 3 et par 9. Il est pair, mais ne finit ni par 0 ni par 5.\end{exemple}
\begin{savoirfaire}[Relier les contraintes]Pour former des lots identiques sans reste, le nombre de lots doit diviser chacun des effectifs. Lister les diviseurs puis garder ceux qui conviennent à toutes les contraintes. Pour des rendez-vous périodiques, comparer les listes de multiples.\end{savoirfaire}
\begin{definition}Prolongement facultatif : un nombre premier a exactement deux diviseurs positifs, 1 et lui-même. 2, 3, 5 et 7 sont premiers ; 1 ne l’est pas. Le crible consiste à barrer les multiples des nombres premiers déjà trouvés.\end{definition}

\clearpage\rubrique{Découverte}
\begin{blocLivre}{ex-2-1}\begin{exercice}[1]{Quotient et reste}
Effectuer la division euclidienne de 67 par 8.
\end{exercice}\end{blocLivre}
\begin{blocLivre}{ex-2-2}\begin{exercice}[1]{Diviseur ou multiple}
Compléter : 54 est un … de 6 ; 6 est un … de 54. Justifier.
\end{exercice}\end{blocLivre}
\begin{blocLivre}{ex-2-3}\begin{exercice}[1]{Tester 450}
450 est-il divisible par 2, 3, 5, 9 et 10 ?
\end{exercice}\end{blocLivre}
\begin{blocLivre}{ex-2-4}\begin{exercice}[1]{Lister des diviseurs}
Donner les diviseurs positifs de 18.
\end{exercice}\end{blocLivre}
\begin{blocLivre}{ex-2-5}\begin{exercice}[1]{Valider une écriture}
$59=6\times8+11$ décrit-il la division euclidienne de 59 par 6 ?
\end{exercice}\end{blocLivre}

\rubrique{Entraînement}
\begin{blocLivre}{ex-2-6}\begin{exercice}[2]{Un chiffre manquant}
Trouver tous les chiffres $a$ tels que le nombre $23a$ soit divisible par 3.
\end{exercice}\end{blocLivre}
\begin{blocLivre}{ex-2-7}\begin{exercice}[2]{Deux critères ensemble}
Trouver le chiffre $a$ pour que $7a5$ soit divisible par 9. Est-il aussi divisible par 5 ?
\end{exercice}\end{blocLivre}
\begin{blocLivre}{ex-2-8}\begin{exercice}[2]{Diviseurs communs}
Lister les diviseurs communs à 24 et 36.
\end{exercice}\end{blocLivre}
\begin{blocLivre}{ex-2-9}\begin{exercice}[2]{Premier multiple commun}
Quel est le plus petit entier strictement positif multiple de 6 et de 8 ?
\end{exercice}\end{blocLivre}
\begin{blocLivre}{ex-2-10}\begin{exercice}[2]{Un crible limité}
Prolongement : trouver les nombres premiers inférieurs à 20.
\end{exercice}\end{blocLivre}

\rubrique{Problèmes}
\begin{blocLivre}{ex-2-11}\begin{exercice}[3]{Transporter une classe}
On transporte 137 personnes dans des minibus de neuf places. Combien faut-il de minibus ? Combien de places resteront libres ?
\end{exercice}\end{blocLivre}
\begin{blocLivre}{ex-2-12}\begin{exercice}[3]{Des bouquets identiques}
On a 42 tulipes et 56 roses. On veut des bouquets identiques sans reste, au moins trois bouquets. Donner tous les nombres de bouquets possibles et la composition du plus grand nombre.
\end{exercice}\end{blocLivre}
\begin{blocLivre}{ex-2-13}\begin{exercice}[3]{Deux signaux}
Un signal se déclenche toutes les 12 minutes et un autre toutes les 18 minutes. Ils sonnent ensemble à 9 h. Quand sonneront-ils ensemble pour la première fois ensuite ?
\end{exercice}\end{blocLivre}
\begin{blocLivre}{ex-2-14}\begin{exercice}[3]{Un nombre à retrouver}
Un entier entre 40 et 60 laisse un reste de 2 quand on le divise par 7 et est divisible par 3. Le trouver.
\end{exercice}\end{blocLivre}
\begin{blocLivre}{ex-2-15}\begin{exercice}[3]{Une affirmation générale}
Mila dit : « La somme de deux multiples de 4 est encore un multiple de 4. » Expliquer avec des paquets, puis vérifier sur 28 et 44.
\end{exercice}\end{blocLivre}

\clearpage
